\makereporttitle{2024}{iunie}

\makereportpreamble

\textbf{Activitate de cercetare:}

În această lună am studiat câteva articole conexe cu cercetarea mea pe zona de haos cuantic, care prezintă niște modele matematice alternative pentru studiul conceptului de haos/complexitate:

\begin{itemize}
    \item „\emph{A Universal Operator Growth Hypothesis}”, de Daniel E. Parker, Xiangyu Cao, Alexander Avdoshkin, Thomas Scaffidi, și Ehud Altman, lucrare publicată în 2024 în Phys. Rev. X 9, 041017.

    \item „\emph{Spectral and Krylov complexity in billiard systems}”, de Hugo A. Camargo, Viktor Jahnke, Hyun-Sik Jeong, Keun-Young Kim, și Mitsuhiro Nishida, lucrare publicată în 2024 în Phys. Rev. D, 109, 046017.

    \item „\emph{Quantum Lyapunov Spectrum}”, de Hrant Gharibyana, Masanori Hanadab, Brian Swinglec și Masaki Tezukad, lucrare publicată în 2019 în Journal of High Energy Physics, Volume 2019, article number 82.
\end{itemize}

\vspace{0.5em}

\textbf{Observațiile conducătorului științific}:

\vspace{3em}

\makereportpostamble{28 iunie 2024}
